\documentclass[11pt]{article}
\usepackage[utf8]{inputenc}
\usepackage{amsmath, amssymb}
\usepackage{geometry}
\usepackage{listings}
\usepackage{xcolor}
\usepackage{graphicx}
\usepackage{tabularx}
\usepackage{booktabs}
\usepackage{hyperref}

\geometry{a4paper, margin=1in}

\definecolor{codegray}{rgb}{0.5,0.5,0.5}
\definecolor{codepurple}{rgb}{0.58,0,0.82}
\definecolor{backcolour}{rgb}{0.95,0.95,0.92}

\lstset{
    language=Python,
    backgroundcolor=\color{backcolour},
    commentstyle=\color{codegray},
    keywordstyle=\color{magenta},
    stringstyle=\color{codepurple},
    basicstyle=\ttfamily\small,
    breaklines=true,
    showstringspaces=false
}

\title{Probabilistic Algorithms: Document Similarity \\ \large Problem Set \& Implementation Guide}
\author{VB}
\date{\today}

\begin{document}

\maketitle

\section*{Problem 0: Start the engine}

Consider the files given.
\begin{center}
\renewcommand{\arraystretch}{1.5}
\begin{tabularx}{\textwidth}{@{} >{\bfseries\raggedright\arraybackslash}p{5.5cm} X @{}}
\toprule
\textcolor{blue}{Filename} & \textcolor{blue}{Purpose} \\ \midrule
\multicolumn{2}{@{}l}{\textbf{Direct Interaction}} \\ \midrule
\lstinline|implementation_tasks.py| & \textcolor{red}{Here you write your code.} Contains the core algorithms for shingling, MinHashing, Bloom Filters, and LSH. \\
\lstinline|lab_dashboard.py| & \textcolor{red}{This file you run.} A testing sandbox with an interactive left sidebar to test each level of your code. \\
\bottomrule
\end{tabularx}
\end{center}

\textbf{Prerequisites:} You will need to install the dependencies for this lab. Run \lstinline|pip install -r requirements.txt| before starting.

Run the file \lstinline|lab_dashboard.py|. This will open the dashboard where you can sequentially test your implementations. You can add your own \lstinline|unittest.TestCase| classes at the bottom of \lstinline|implementation_tasks.py| and run \lstinline|python implementation_tasks.py| to execute them.

\subsubsection*{Self-Testing (Optional)}
At the bottom of \lstinline|implementation_tasks.py|, you will find a block that looks like this:
\begin{lstlisting}[language=Python]
if __name__ == '__main__':
    import unittest
    # Add your own unittest.TestCase classes here...
    unittest.main(verbosity=2)
\end{lstlisting}
This allows you to write your own automated unit tests to verify your code before running the visual dashboard. If you'd like to test your functions with specific inputs, simply define a class inheriting from \lstinline|unittest.TestCase| above the \lstinline|unittest.main()| call. For example, if you wanted to test simple math:
\begin{lstlisting}[language=Python]
class TestMyMath(unittest.TestCase):
    def test_addition(self):
        result = 2 + 2
        self.assertEqual(result, 4, "Addition is broken!")
\end{lstlisting}
You can run these tests anytime by executing \lstinline|python implementation_tasks.py| in your terminal. This is completely optional, but highly recommended for debugging tricky edge cases.

\section*{Problem 1: K-Shingling (Level 1)}

\subsubsection*{Mathematical part}
To compare two text documents, we first need to convert them into a mathematical set. We do this using \textbf{k-shingling}. A $k$-shingle is any continuous sequence of $k$ words (or characters) appearing in the document.

\textbf{Example:} Let the text be ``The quick brown fox jumps''. For $k=3$, the shingles are:
$$ \{ (\text{the, quick, brown}), (\text{quick, brown, fox}), (\text{brown, fox, jumps}) \} $$

\subsubsection*{Implementation part}
\textbf{Implement:} \lstinline|get_shingles(text, k=3)|.
1. Convert the text to lowercase.
2. Extract all words (e.g., using a regular expression like \lstinline|r'\b\w+\b'|).
3. Create tuples of length $k$ by sliding a window over the words.
4. Return a Python \lstinline|set| containing these tuples.


\section*{Problem 2: Jaccard Similarity (Level 2)}

\subsubsection*{Mathematical part}
The \textbf{Jaccard Similarity} between two sets $A$ and $B$ is the size of their intersection divided by the size of their union:
$$ J(A, B) = \frac{|A \cap B|}{|A \cup B|} $$
This gives a value between 0 (completely disjoint) and 1 (identical sets).

\subsubsection*{Implementation part}
\textbf{Implement:} \lstinline|jaccard_similarity(set1, set2)|. Return the float value. Be sure to handle the edge case where both sets are empty (return 0.0).


\section*{Problem 3: MinHash Signatures (Level 3)}

\subsubsection*{Mathematical part}
If we have 10,000 documents, doing a pairwise Jaccard comparison takes $O(N^2)$ operations on massive sets, which is too slow.
Instead, we use \textbf{MinHash}. If we pass every element of a set through a random hash function $h(x)$, the probability that the minimum hash value of set $A$ equals the minimum hash value of set $B$ is exactly their Jaccard similarity!
$$ P(\min(h(A)) == \min(h(B))) = J(A, B) $$
By doing this for 100 different hash functions, we can compress a massive document into a tiny ``signature'' of 100 integers.

\subsubsection*{Implementation part}
\textbf{Implement:} \lstinline|create_signature(shingles, hash_funcs)|.
For every hash function in the provided list \lstinline|hash_funcs|, find the minimum hash value it produces across all shingles in the document. Return a tuple of these minimums.


\section*{Problem 4: MinHash Similarity (Level 4)}

\subsubsection*{Mathematical part}
Once we have our MinHash signatures, how do we compare them? Because $P(\min(h(A)) == \min(h(B))) = J(A, B)$, we can estimate the Jaccard similarity by simply counting the fraction of hash values that are identical between the two signatures!
$$ \hat{J}(A, B) \approx \frac{\text{Number of identical hash values}}{\text{Total number of hash functions}} $$

\subsubsection*{Implementation part}
\textbf{Implement:} \lstinline|minhash_similarity(sig1, sig2)|.
Compare two tuples of integers element-by-element. Count how many indices have the exact same value. Use the \textbf{MinHash Sandbox} tab in the dashboard to test your estimation against exact Jaccard similarity.


\section*{Problem 5: Modified Bloom Filter (LSH Bucketing) (Level 5)}

\subsubsection*{Mathematical part}
Even with tiny signatures, comparing 10,000 signatures pairwise is still $O(N^2)$. We only want to compare documents that are likely to be similar. 

We can think of this process as a \textbf{Modified Bloom Filter} with memory. We divide the signature into $b$ \textbf{bands} of $r$ \textbf{rows}. We hash each band into a bucket. Unlike a standard Bloom Filter that only stores a `1` bit to indicate existence, our bucket "remembers" exactly which documents were hashed into it by storing a list of Document IDs. If two documents hash to the same bucket in any band, they collide!

\subsubsection*{Implementation part}
\textbf{Implement:} \lstinline|lsh_bloom_filter(signatures, num_bands)|.
Iterate over the signatures. For each signature, split it into \lstinline|num_bands|. Use the band index and the tuple of values in that band to create a unique bucket identifier. Return a dictionary mapping these bucket IDs to lists of document indices. Use the \textbf{Bloom Simulator} tab to watch the buckets fill up!


\section*{Problem 6: Standard Bloom Filter Probability (Level 6)}

\subsubsection*{Mathematical part}
Let's step back and look at a Standard Bloom Filter. If we have a filter of $m$ bits, and we insert $n$ elements using $k$ independent hash functions, what is the probability of a False Positive (a collision)?
The probability that a specific bit is still 0 after all insertions is approximately $e^{-kn/m}$. Therefore, the probability that all $k$ bits for a new element are already 1 (a False Positive) is:
$$ P(\text{False Positive}) \approx \left(1 - e^{-k\frac{n}{m}}\right)^k $$

\subsubsection*{Implementation part}
\textbf{Implement:} \lstinline|bloom_false_positive(n, m, k)|.
Return the probability. You can visualize this curve in the \textbf{Bloom Probabilities} tab.


\section*{Problem 7: Document Matcher (Level 7)}

Now that all your core hashing and filtering logic is complete, open the \textbf{Document Matcher} tab. This runs the full pipeline on 300 real NASA press releases. If your functions work, you will instantly see all the duplicated or highly similar articles!


\section*{Problem 8: LSH Probability Mechanics (Level 8)}

\subsubsection*{Mathematical part}
To understand why the Modified Bloom Filter (LSH) is so effective, we must look at the probability curve of a collision. 
If we divide the signature into $b$ bands of $r$ rows, the probability that two documents with Jaccard similarity $s$ will collide in at least one band is:
$$ P(\text{collision}) = 1 - (1 - s^r)^b $$
This creates a sharp ``S-curve''. The threshold similarity (where the probability of collision crosses 50\%) can be approximated as:
$$ t \approx \left(1 - \left(\frac{1}{2}\right)^{\frac{1}{b}}\right)^{\frac{1}{r}} $$

\subsubsection*{Implementation part}
\textbf{Implement:}
1. \lstinline|collision_probability(jaccard, bands, rows)|: Return $1 - (1 - j^r)^b$.
2. \lstinline|calculate_threshold(bands, rows)|: Return $(1 - (1/2)^{1/b})^{1/r}$.
Check the \textbf{LSH Probabilities} tab to see your S-Curve perfectly separate False Positives from False Negatives!

\section*{\textcolor{red}{Submission}}

Submit the file \texttt{implementation\_tasks.py} once the dashboard successfully catches the duplicated press releases!

\end{document}
